% Einheit 1 von 3 – Vektorbegriff, Betrag und Kollinearität (bank/vektoren-und-rechenoperationen/e1.jsonl, zusammenbau v0.1)
%% TODO Zeitmarke Einheit 1: Marken-Zeile nicht lesbar
%% TODO Zweigzeile Teil 1: was hier gelernt wird (Satz in Schülersprache aus dem Einheitstitel) – Einheit 1
\einheitenkopf[e1]{Einheit 1 von 3 · Vektorbegriff, Betrag und Kollinearität}
\zweigzeile{keine Prüfungsaufgabe}
\setcounter{aufgabe}{8}

%% TODO Titel: Ich-kann-Satz fehlt in der Bank (Platzhalter: Kettenname) – Nr. 9
%% TODO Anweisung: ein Satz über der ersten Teilaufgabe fehlt in der Bank – Nr. 9
\begin{aufgabe}{Faktor gesucht oder Richtung gesucht?}
\begin{teile}
\teil Gesucht ist $a$ in $\vec v = (a | 8)$ mit $|\vec v| = 17$. Was fehlt? \\ \kreuz{Zahl} \\ \kreuz{Richtung}
\end{teile}
\end{aufgabe}

%% TODO Titel: Ich-kann-Satz fehlt in der Bank (Platzhalter: Kettenname) – Nr. 10
%% TODO Anweisung: ein Satz über der ersten Teilaufgabe fehlt in der Bank – Nr. 10
\begin{aufgabe}{Vektorbegriff und Betrag}
%% TODO \rechenplatz{n}: Schrittzahl des Grundfalls fehlt in der Bank (nur bei fester Schreibform, 2.3 b)
\begin{teile}
\teil Die Länge der Strecke von $A$ nach $B$: Zahl oder Vektor? \\ \kreuz{Zahl} \\ \kreuz{Vektor}
\teil $A(1 | 2 | 0)$, $B(4 | 3 | 2)$: $\overrightarrow{AB}$, $\overrightarrow{BA}$ und $2 \cdot \overrightarrow{AB}$?
\teil $\vec u = (2 | 3 | 6)$, $\vec v = (4 | 4 | 2)$: welcher ist länger? |u| = \leerfeld , |v| = \leerfeld
\teil $\vec v = (a | 12)$, $a > 0$, $|\vec v| = 13$: $a$? \hfill (Abitur 2025 LK) \\ a = \leerfeld
\teil $A(2 | 0 | 1)$, $B(2 | 3 | 3)$, $D(2 | 1 | 0)$, $C(2 | y | 4)$, die Kante $DC$ ist parallel zu $AB$: $y$? \hfill (Abitur 2025 LK) \\ y = \leerfeld
\teil Zwei Stäbe verbinden $(0 | 0 | 0)$ mit $(4 | 4 | 4)$ und $(0 | 4 | 0)$ mit $(4 | 0 | 4)$. Gib eine Blickrichtung an, aus der man ein X sieht, und zeichne die Ansicht von oben. \hfill (Abitur 2022 LK)

\begin{ksys3} \rgerade[0:1]{0,0,0}{4,4,4}{} \rgerade[0:1]{0,4,0}{4,-4,4}{} \end{ksys3} \begin{ksys}[xmin=0,xmax=5,ymin=0,ymax=5,xlabel=$x_1$,ylabel=$x_2$] \end{ksys}
\teil Ein Verleih verteilt 600 Leihräder auf drei Stationen; jede Verteilung wird als Vektor (Räder an Station 1 | 2 | 3) notiert. Können zwei solche Vektoren verschiedene Beträge haben? Begründe mit einem Beispiel. \hfill (Abitur 2024 GK)
\end{teile}
\end{aufgabe}

%% TODO Titel: Ich-kann-Satz fehlt in der Bank (Platzhalter: Kettenname) – Nr. 11
%% TODO Anweisung: ein Satz über der ersten Teilaufgabe fehlt in der Bank – Nr. 11
\begin{aufgabe}{Vektorbegriff und Betrag – Fehler finden, Begründen, Darstellungswechsel, Anwendung}
\begin{teile}
\teil Ich finde den Fehler: Mia rechnet $|(4 | 4 | 7)| = 4 + 4 + 7 = 15$.
\teil Warum liefert $a^2 = 25$ zwei Werte für $a$?
\teil Lies $\overrightarrow{AB}$ ab.

\begin{ksys}[xmin=-1,xmax=5,ymin=-1,ymax=5,ablesen] \punkt{1}{1}{A} \punkt{4}{3}{B} \end{ksys}
( \leerfeld | \leerfeld )
\teil Eine Drohne fliegt geradlinig von $S(0 | 0 | 2)$ nach $Z(30 | 40 | 2)$ (Angaben in m). Der Akku reicht noch für 45 m. Kommt sie an?
\end{teile}
\end{aufgabe}
